% Chapter 6: Lieb's Concavity and Ando's Convexity
\chapter{Lieb's Concavity and Ando's Convexity}

This chapter contains the main results of the project.
Working in an abstract H$^*$-algebra $H$, we derive Lieb's concavity theorem and
Ando's convexity theorem by specializing the operator power mean from
Chapter~5 to the pair of bounded operators $(R_B, L_A)$ acting on $H$,
and then reading off the resulting joint concavity or convexity of the
functional $T \mapsto \operatorname{re}\langle x,\, T x\rangle$.

\section{The Power Mean Evaluated at a Vector}

\begin{proposition}[Simplified formula for the operator power mean]
  \label{prop:powermean_apply}
  \lean{PowerMean_apply}
  \leanok
  \uses{thm:perspective_rpow_commute, prop:mul_operators_cfc}
  Let $H$ be an H$^*$-algebra, $\alpha, \beta \in \mathbb{R}$ with $\alpha \geq 0$ and
  $\beta \neq 0$.  Let $A$ be strictly positive and $B \geq 0$.  Then for all $x \in H$,
  \[
    \Pi_{\alpha,\beta}(L_A, R_B)\,x \;=\; A^{\beta(1-\alpha)}\, x\, B^\alpha.
  \]
  \begin{proof}
    \proves{prop:powermean_apply}
    \leanok
    \uses{thm:perspective_rpow_commute, prop:mul_operators_cfc}
    The power mean is defined as $\Pi_{\alpha,\beta}(L_A, R_B) =
    ((\cdot)^\alpha \triangle (\cdot)^\beta)(R_B, L_A)$.
    Since $L_A$ and $R_B$ commute and $R_B \geq 0$, $L_A$ is strictly positive,
    the generalized perspective simplifies by Theorem~\ref{thm:perspective_rpow_commute}
    (which requires Lemma~\ref{lem:mul_rpow_commute}), giving
    $(R_B)^\alpha \cdot (L_A)^{\beta(1-\alpha)}$.
    Evaluating at $x$ and using $L_A^r\,x = A^r x$ and $R_B^r\,x = x\,B^r$
    yields $A^{\beta(1-\alpha)}\,x\,B^\alpha$.
  \end{proof}
\end{proposition}

\section{General H\texorpdfstring{$^*$}{*}-Algebra Versions}

\begin{theorem}[Lieb concavity --- general form]
  \label{thm:lieb_general}
  \lean{LiebConcavity_general}
  \leanok
  \uses{def:hstar_algebra}
  Let $H$ be an H$^*$-algebra, $\alpha, \beta \in (0,1]$, and $x \in H$.
  Then the map
  \[
    (A, B) \;\longmapsto\; \operatorname{re}\bigl\langle x,\;
      A^{\beta(1-\alpha)}\, x\, B^{\alpha} \bigr\rangle
  \]
  is jointly concave on
  $\{(A, B) \in H \times H \mid A \text{ strictly positive},\, B \geq 0\}$.
  \begin{proof}
    \proves{thm:lieb_general}
    \leanok
    \uses{thm:power_mean_concave, prop:powermean_apply, lem:reapply_inner_mono}
    Consider the operator power mean $\Pi_{\alpha,\beta}(L_A, R_B)$ acting on
    $H \to_L H$, which by definition is the generalized perspective of the
    $\alpha$-power mean at the pair $(L_A, R_B)$.
    By the CFC commutativity of left and right multiplication
    (Proposition~\ref{prop:mul_operators_cfc}), we have
    $L_{A^{\beta(1-\alpha)}} = (L_A)^{\beta(1-\alpha)}$ and
    $R_{B^\alpha} = (R_B)^\alpha$, so the power mean evaluated at $x$ equals
    $A^{\beta(1-\alpha)}\, x\, B^\alpha$.
    Joint concavity of $(A, B) \mapsto \Pi_{\alpha,\beta}(L_A, R_B)$ in the
    operator order on $H \to_L H$ then follows from
    Theorem~\ref{thm:power_mean_concave}.
    Finally, the map $T \mapsto \operatorname{re}\langle x, T x\rangle$ is linear
    and hence preserves concavity, yielding the stated result.
  \end{proof}
\end{theorem}

\begin{theorem}[Ando convexity --- general form]
  \label{thm:ando_general}
  \lean{AndoConvexity_general}
  \leanok
  \uses{def:hstar_algebra}
  Let $H$ be an H$^*$-algebra, $\alpha \in [1,2]$, $\beta \in (0,1]$, and $x \in H$.
  Then the map
  \[
    (A, B) \;\longmapsto\; \operatorname{re}\bigl\langle x,\;
      A^{\beta(1-\alpha)}\, x\, B^{\alpha} \bigr\rangle
  \]
  is jointly convex on
  $\{(A, B) \in H \times H \mid A \text{ strictly positive},\, B \geq 0\}$.
  \begin{proof}
    \proves{thm:ando_general}
    \leanok
    \uses{thm:power_mean_convex, prop:powermean_apply, lem:reapply_inner_mono}
    The argument is analogous to Theorem~\ref{thm:lieb_general}.
    The same CFC commutativity identities (Proposition~\ref{prop:mul_operators_cfc})
    identify the power mean evaluated at $x$ with
    $A^{\beta(1-\alpha)}\, x\, B^\alpha$.
    Since $\alpha \in [1,2]$, the power mean is jointly convex in operator order
    by Theorem~\ref{thm:power_mean_convex}, and linearity of
    $T \mapsto \operatorname{re}\langle x, T x\rangle$ propagates convexity
    to the scalar functional.
  \end{proof}
\end{theorem}

\section{Lieb's Extension Theorem}

\begin{theorem}[Lieb's Extension Theorem]
  \label{thm:lieb_extension}
  \lean{LiebExtension}
  \leanok
  \uses{def:hstar_algebra}
  Let $H$ be an H$^*$-algebra, $p, q > 0$ with $p + q \leq 1$, and $x \in H$.
  Then the map
  \[
    (A, B) \;\longmapsto\; \operatorname{re}\bigl\langle x,\;
      A^{q}\, x\, B^{p} \bigr\rangle
  \]
  is jointly concave on
  $\{(A, B) \in H \times H \mid A \text{ strictly positive},\, B \geq 0\}$.
  \begin{proof}
    \proves{thm:lieb_extension}
    \leanok
    \uses{thm:lieb_general}
    Set $\alpha := p$ and $\beta := q / (1 - p)$.
    The hypothesis $p + q \leq 1$ with $p, q > 0$ ensures $\beta \leq 1$,
    so $(\alpha, \beta) \in (0,1] \times (0,1]$, and one computes
    $\beta(1-\alpha) = q$.
    The result therefore follows from Theorem~\ref{thm:lieb_general} with these
    parameter choices.
  \end{proof}
\end{theorem}

\section{Lieb's Concavity Theorem}

\begin{corollary}[Lieb's Concavity Theorem]
  \label{thm:lieb}
  \lean{LiebConcavity}
  \leanok
  \uses{def:hstar_algebra}
  Let $H$ be an H$^*$-algebra, $s \in (0,1)$, and $x \in H$.
  Then the map
  \[
    (A, B) \;\longmapsto\; \operatorname{re}\bigl\langle x,\;
      A^{s}\, x\, B^{1-s} \bigr\rangle
  \]
  is jointly concave on
  $\{(A, B) \in H \times H \mid A \text{ strictly positive},\, B \geq 0\}$.
  \begin{proof}
    \proves{thm:lieb}
    \leanok
    \uses{thm:lieb_extension}
    Apply Theorem~\ref{thm:lieb_extension} with $q := s$ and $p := 1 - s$,
    noting that $p, q > 0$ and $p + q = 1 \leq 1$.
  \end{proof}
\end{corollary}

\section{Ando's Convexity Theorem}

\begin{theorem}[Ando's Convexity Theorem]
  \label{thm:ando}
  \lean{AndoConvexity}
  \leanok
  \uses{def:hstar_algebra}
  Let $H$ be an H$^*$-algebra, $q \in [1,2]$, $r > 0$ with $q - r > 1$,
  and $x \in H$.
  Then the map
  \[
    (A, B) \;\longmapsto\; \operatorname{re}\bigl\langle x,\;
      A^{-r}\, x\, B^{q} \bigr\rangle
  \]
  is jointly convex on
  $\{(A, B) \in H \times H \mid A \text{ strictly positive},\, B \geq 0\}$.
  \begin{proof}
    \proves{thm:ando}
    \leanok
    \uses{thm:ando_general}
    Set $\alpha := q$ and $\beta := r / (q - 1)$.
    The hypothesis $q - r > 1$ with $r > 0$ and $q \in [1,2]$ ensures
    $\beta \in (0,1]$, and one computes $\beta(1-\alpha) = -r$.
    The result follows from Theorem~\ref{thm:ando_general} with these
    parameter choices.
  \end{proof}
\end{theorem}
